\subsection{Backend Architecture}
\label{subsec:backend-architecture}

The backend infrastructure of the speech training application implements a modular, scalable architecture based on Clean Architecture principles \cite{martin2017clean}. The system is designed to facilitate real-time speech analysis through AI-powered inference while maintaining strict separation of concerns and testability.

\subsubsection{Technology Stack}

The backend is implemented using modern, production-ready technologies optimized for high-performance real-time processing:

\begin{itemize}
    \item \textbf{Framework:} FastAPI (Python 3.13) -- asynchronous REST API framework with automatic OpenAPI documentation and native WebSocket support
    \item \textbf{Database:} PostgreSQL 17 with TimescaleDB extension -- providing time-series optimization for session and performance tracking data
    \item \textbf{Containerization:} Docker Compose -- ensuring consistent development and deployment environments with hot-reload capabilities
    \item \textbf{AI/ML Frameworks:} PyTorch and Transformers library -- for custom phoneme recognition models and automatic speech recognition (ASR)
    \item \textbf{Audio Processing:} Faster-Whisper -- optimized OpenAI Whisper implementation for real-time speech-to-text transcription
\end{itemize}

\subsubsection{Architectural Design}

The backend follows the Clean Architecture (Hexagonal Architecture) pattern \cite{cockburn2005hexagonal}, enforcing strict dependency inversion and separation of business logic from infrastructure concerns. This design ensures maintainability, testability, and adaptability to changing requirements.

The architecture is organized into four distinct layers with unidirectional dependencies flowing inward toward the domain core (Figure \ref{fig:backend-architecture}):

\begin{description}
    \item[Domain Layer (Core):] Contains enterprise business rules, domain entities (e.g., \texttt{User}, \texttt{Exercise}, \texttt{PracticeSession}), value objects, and port interfaces defining contracts for external dependencies. This layer has zero external dependencies.
    
    \item[Application Layer:] Implements use case orchestration and application-specific business rules through service classes. Services coordinate domain entities and interact with infrastructure through port interfaces defined in the domain layer.
    
    \item[Infrastructure Layer:] Provides concrete implementations of domain ports, including:
    \begin{itemize}
        \item Database persistence via SQLAlchemy ORM repositories
        \item AI model adapters for phoneme recognition and ASR
        \item Audio preprocessing components
        \item External service integrations
    \end{itemize}
    
    \item[API Layer (Presentation):] Exposes REST endpoints and WebSocket connections through FastAPI routers. Request/response transformation is handled by Pydantic schemas, with dependency injection managing service instantiation.
\end{description}

\subsubsection{AI Model Integration}

The system integrates custom-trained transformer-based models for speech analysis through a flexible adapter pattern. The domain layer defines an \texttt{IASRService} interface, with concrete implementations in the infrastructure layer:

\begin{itemize}
    \item \textbf{Phoneme Recognition:} Custom Transformer-CTC models trained for German phoneme sequences, enabling fine-grained pronunciation analysis
    \item \textbf{Automatic Speech Recognition:} Faster-Whisper adapter for OpenAI Whisper models, providing word-level transcription with timestamps
    \item \textbf{Audio Preprocessing:} Standardized pipeline for audio normalization, resampling, and feature extraction
\end{itemize}

Models are loaded at application startup and accessed through dependency injection, ensuring efficient resource utilization and simplified testing through mock implementations.

\subsubsection{Real-Time WebSocket Communication}

Speech analysis requires real-time progress feedback during long-running inference tasks. The system implements a WebSocket-based communication layer for bidirectional client-server interaction:

\begin{enumerate}
    \item Client initiates analysis via REST POST request to \texttt{/api/v1/analyze/sessions/\{id\}}
    \item Server responds with WebSocket connection URL and immediately starts background processing
    \item Client establishes WebSocket connection to \texttt{/api/v1/analyze/ws/sessions/\{id\}}
    \item Server streams progress updates (percentage, current item, intermediate results) in JSON format
    \item Upon completion, server sends final results and closes connection
\end{enumerate}

The \texttt{ConnectionManager} component handles WebSocket lifecycle, supporting group-based broadcasting for multi-user scenarios and automatic cleanup on disconnection. Authentication is enforced through JWT tokens passed in the WebSocket handshake headers.

\begin{figure*}[t]
\centering
\begin{adjustbox}{max width=\textwidth}
\begin{tikzpicture}[
    node distance=1.2cm,
    layer/.style={rectangle, draw, thick, minimum width=3.5cm, minimum height=1cm, align=center, fill=blue!10},
    component/.style={rectangle, draw, minimum width=3cm, minimum height=0.7cm, align=center, fill=white},
    arrow/.style={->, thick, >=stealth},
    label/.style={font=\small\bfseries}
]

% Layers (from top to bottom)
\node[layer, minimum width=12cm] (api) {\textbf{API Layer (Presentation)}\\FastAPI Routers, WebSocket Handlers};
\node[layer, minimum width=12cm, below=of api] (app) {\textbf{Application Layer}\\Use Case Services, Business Logic};
\node[layer, minimum width=12cm, below=of app] (core) {\textbf{Domain Layer (Core)}\\Entities, Value Objects, Port Interfaces};
\node[layer, minimum width=12cm, below=of core] (infra) {\textbf{Infrastructure Layer}\\Repository Implementations, AI Model Adapters};

% Infrastructure components (below infra layer)
\node[component, below left=0.8cm and -2cm of infra] (db) {PostgreSQL\\TimescaleDB};
\node[component, below=0.8cm of infra] (models) {AI Models\\Phoneme/ASR};
\node[component, below right=0.8cm and -2cm of infra] (storage) {File Storage\\Media};

% External clients (above API layer)
\node[component, above=0.8cm of api] (client) {Web Client\\Vue.js};

% Arrows - Dependency flow
\draw[arrow] (api) -- (app) node[midway, right, font=\footnotesize] {uses};
\draw[arrow] (app) -- (core) node[midway, right, font=\footnotesize] {uses};
\draw[arrow] (infra) -- (core) node[midway, right, font=\footnotesize] {implements};
\draw[arrow] (api) -- (core);

% Infrastructure connections
\draw[arrow] (infra) -- (db);
\draw[arrow] (infra) -- (models);
\draw[arrow] (infra) -- (storage);

% Client connections
\draw[arrow, <->, red, thick] (client) -- (api) node[midway, right, font=\footnotesize, text=red] {REST/WS};

% Side annotations
\node[label, right=2.5cm of app, text width=3.5cm] (di) {Dependency\\Injection\\Container};
\node[label, right=2.5cm of core, text width=3.5cm] (ports) {Port Interfaces\\(Abstractions)};

% Dependency injection arrow
\draw[arrow, dashed, blue] (di) -- (app);
\draw[arrow, dashed, blue] (di) -- (api);

% Key components detail (right side)
\node[component, right=5.5cm of api, minimum width=4cm, fill=red!10] (ws) {WebSocket\\Connection Manager};
\node[component, below=0.5cm of ws, minimum width=4cm, fill=green!10] (jwt) {JWT Authentication\\Middleware};
\node[component, below=0.5cm of jwt, minimum width=4cm, fill=yellow!10] (asr) {ASR Service\\Interface};

\draw[arrow, dashed] (api) -- (ws);
\draw[arrow, dashed] (api) -- (jwt);
\draw[arrow, dashed] (infra) -- (asr);

% Legend
\node[below=2.5cm of db, font=\footnotesize, text width=11cm, align=left] (legend) {
    \textbf{Dependency Direction:} Infrastructure/API $\rightarrow$ Application $\rightarrow$ Domain Core\\
    \textbf{Communication:} \textcolor{red}{Synchronous REST \& Asynchronous WebSocket}
};

\end{tikzpicture}
\end{adjustbox}
\caption{Backend architecture following Clean Architecture principles with layer separation and dependency inversion. The WebSocket connection manager enables real-time AI inference progress updates.}
\label{fig:backend-architecture}
\end{figure*}

\subsubsection{Deployment and Development Workflow}

The entire backend infrastructure runs within Docker containers, ensuring environment consistency across development, testing, and production. Key deployment characteristics include:

\begin{itemize}
    \item \textbf{Hot-reload:} Source code changes trigger automatic service restart without container rebuilds
    \item \textbf{Database Migrations:} Alembic manages schema versioning with automatic migration generation
    \item \textbf{Testing:} Isolated test database instances with pytest-based unit, integration, and performance tests
    \item \textbf{API Documentation:} Automatic OpenAPI/Swagger documentation at \texttt{/docs} endpoint
\end{itemize}

All backend operations (migrations, tests, shell access) execute within the containerized environment using \texttt{docker compose exec}, eliminating dependency conflicts and ensuring reproducible builds.
